\documentclass[11pt,a4paper]{article}

\usepackage[utf8]{inputenc}
\usepackage[T1]{fontenc}
\usepackage[italian]{babel}
\usepackage{amsmath,amssymb,amsthm}
\usepackage{algorithm}
\usepackage[noend]{algpseudocode}
\usepackage{booktabs}
\usepackage{geometry}
\usepackage{hyperref}
\geometry{margin=2.5cm}

\newcommand{\bx}{\mathbf{x}}
\newcommand{\bmu}{\boldsymbol{\mu}}
\newcommand{\bth}{\boldsymbol{\theta}}
\newcommand{\Nrm}{\mathcal{N}}
\newcommand{\E}{\mathbb{E}}

\theoremstyle{definition}
\newtheorem{prop}{Proposizione}

\title{Expectation-Maximization per Gaussian Mixture Model:\\
base teorica e fonti di parallelismo}
\author{Parallel Expectation-Maximization Clustering -- HPC4DS}
\date{}

\begin{document}
\maketitle

\begin{abstract}
Questo documento riassume la teoria dell'algoritmo Expectation-Maximization (EM)
applicato ai Gaussian Mixture Model (GMM) con covarianza diagonale, e spiega
perché l'algoritmo si presta a una parallelizzazione efficiente con MPI e
OpenMP. L'idea centrale è che tutte le operazioni costose sono \emph{somme su
punti indipendenti}: i dati si partizionano, i parametri (molto piccoli) si
replicano e si comunicano solo somme parziali.
\end{abstract}

%------------------------------------------------------------------
\section{Il problema}
%------------------------------------------------------------------
Dato un insieme di $N$ punti $X=\{\bx_1,\dots,\bx_N\}$, con $\bx_i\in\mathbb{R}^D$,
vogliamo raggrupparli in $K$ cluster in modo \emph{soft}: ogni punto ha una
probabilità di appartenere a ciascun cluster, non un'etichetta rigida come in
$k$-means.

\paragraph{Notazione.}
\begin{center}
\begin{tabular}{ll}
\toprule
Simbolo & Significato \\
\midrule
$N$ & numero di punti \\
$D$ & numero di feature \\
$K$ & numero di cluster (componenti) \\
$\pi_k$ & peso della componente $k$, $\sum_k\pi_k=1$, $\pi_k\ge 0$ \\
$\bmu_k\in\mathbb{R}^D$ & media della componente $k$ \\
$\sigma_{kd}^2$ & varianza della componente $k$ sulla feature $d$ \\
$\bth$ & insieme dei parametri $\{\pi_k,\bmu_k,\sigma^2_{k\cdot}\}_{k=1}^K$ \\
$z_i\in\{1,\dots,K\}$ & variabile latente: cluster che ha generato $\bx_i$ \\
$r_{ik}$ & responsabilità: $P(z_i=k\mid\bx_i,\bth)$ \\
\bottomrule
\end{tabular}
\end{center}

%------------------------------------------------------------------
\section{Gaussian Mixture Model}
%------------------------------------------------------------------
Si assume che ogni punto sia generato così: si estrae un cluster
$z_i\sim\mathrm{Categorical}(\pi_1,\dots,\pi_K)$, poi
$\bx_i\mid z_i=k\sim\Nrm(\bmu_k,\Sigma_k)$. La densità marginale è
\begin{equation}
p(\bx_i\mid\bth)=\sum_{k=1}^{K}\pi_k\,\Nrm(\bx_i\mid\bmu_k,\Sigma_k).
\end{equation}
Con \textbf{covarianza diagonale}, $\Sigma_k=\mathrm{diag}(\sigma_{k1}^2,\dots,\sigma_{kD}^2)$, la
gaussiana si fattorizza sulle feature:
\begin{equation}
\Nrm(\bx\mid\bmu_k,\Sigma_k)=\prod_{d=1}^{D}\frac{1}{\sqrt{2\pi\sigma_{kd}^2}}
\exp\!\left(-\frac{(x_d-\mu_{kd})^2}{2\sigma_{kd}^2}\right).
\label{eq:gauss}
\end{equation}
Il costo di una valutazione è $O(D)$ (con covarianza piena sarebbe $O(D^2)$ o
$O(D^3)$).

\subsection{Log-likelihood}
I punti sono assunti indipendenti, quindi
\begin{equation}
\ell(\bth)=\sum_{i=1}^{N}\log\sum_{k=1}^{K}\pi_k\,\Nrm(\bx_i\mid\bmu_k,\Sigma_k).
\label{eq:loglik}
\end{equation}
Massimizzare $\ell$ rispetto a $\bth$ direttamente è difficile: il $\log$ della
somma interna impedisce una soluzione in forma chiusa. Se conoscessimo gli $z_i$
il problema sarebbe banale (stime di massima verosimiglianza per ogni cluster
separatamente). EM sfrutta esattamente questo.

%------------------------------------------------------------------
\section{L'algoritmo EM}
%------------------------------------------------------------------
EM alterna due passi, partendo da un'inizializzazione $\bth^{(0)}$.

\subsection{E-step: calcolo delle responsabilità}
Con i parametri correnti $\bth^{(t)}$, la probabilità a posteriori che $\bx_i$
provenga dal cluster $k$ è, per la regola di Bayes,
\begin{equation}
r_{ik}=P(z_i=k\mid\bx_i,\bth^{(t)})=
\frac{\pi_k\,\Nrm(\bx_i\mid\bmu_k,\Sigma_k)}
{\sum_{j=1}^{K}\pi_j\,\Nrm(\bx_i\mid\bmu_j,\Sigma_j)}.
\label{eq:estep}
\end{equation}
Per ogni $i$, $\sum_k r_{ik}=1$.

\subsection{M-step: aggiornamento dei parametri}
Si massimizza il valore atteso della log-likelihood completa
$Q(\bth\mid\bth^{(t)})=\E_{z\mid X,\bth^{(t)}}[\log p(X,z\mid\bth)]$. Con le
responsabilità fissate si ottengono formule chiuse. Definiamo le
\emph{statistiche sufficienti}
\begin{align}
N_k &= \sum_{i=1}^{N} r_{ik}, &
S_{kd}^{(1)} &= \sum_{i=1}^{N} r_{ik}\,x_{id}, &
S_{kd}^{(2)} &= \sum_{i=1}^{N} r_{ik}\,x_{id}^2.
\label{eq:suff}
\end{align}
I nuovi parametri sono
\begin{align}
\pi_k^{(t+1)} &= \frac{N_k}{N}, \label{eq:pi}\\
\mu_{kd}^{(t+1)} &= \frac{S_{kd}^{(1)}}{N_k}, \label{eq:mu}\\
\sigma_{kd}^{2\,(t+1)} &= \frac{S_{kd}^{(2)}}{N_k}-\left(\mu_{kd}^{(t+1)}\right)^2.
\label{eq:var}
\end{align}
La forma \eqref{eq:var} usa $\E[x^2]-\E[x]^2$ ed è comoda perché richiede solo
somme che si possono accumulare in un solo passaggio sui dati (nella pratica si
aggiunge un floor $\sigma^2\ge\varepsilon$ per evitare varianze nulle).

\subsection{Convergenza}
Si calcola $\ell(\bth^{(t)})$ con \eqref{eq:loglik} e ci si ferma quando
$|\ell^{(t+1)}-\ell^{(t)}|<\varepsilon$ oppure si raggiunge un numero massimo di
iterazioni.

\begin{prop}[Monotonia]
Ad ogni iterazione EM, $\ell(\bth^{(t+1)})\ge\ell(\bth^{(t)})$.
\end{prop}
\noindent
Quindi EM converge a un punto stazionario, in generale un massimo \emph{locale}:
il risultato dipende dall'inizializzazione. Per confrontare versioni diverse
(sequenziale, MPI, ottimizzazioni) l'inizializzazione deve essere identica.

\begin{algorithm}[h]
\caption{EM per GMM a covarianza diagonale (sequenziale)}
\begin{algorithmic}[1]
\State Inizializza $\bth^{(0)}$; $t\gets 0$
\Repeat
  \For{$i=1$ \textbf{to} $N$} \Comment{E-step}
    \For{$k=1$ \textbf{to} $K$}
      \State $a_{ik}\gets\log\pi_k+\log\Nrm(\bx_i\mid\bmu_k,\Sigma_k)$ \Comment{$O(D)$}
    \EndFor
    \State $m_i\gets\log\sum_k e^{a_{ik}}$ \Comment{log-sum-exp}
    \State $r_{ik}\gets e^{a_{ik}-m_i}$ per ogni $k$
  \EndFor
  \State $\ell\gets\sum_i m_i$
  \State calcola $N_k,S^{(1)}_{kd},S^{(2)}_{kd}$ con \eqref{eq:suff} \Comment{M-step}
  \State aggiorna $\pi,\bmu,\sigma^2$ con \eqref{eq:pi}--\eqref{eq:var}
  \State $t\gets t+1$
\Until{$|\ell^{(t)}-\ell^{(t-1)}|<\varepsilon$}
\end{algorithmic}
\end{algorithm}

\paragraph{Stabilità numerica.}
Le densità \eqref{eq:gauss} possono andare in underflow quando $D$ è grande. Si
lavora in scala logaritmica e si usa il trucco log-sum-exp:
$\log\sum_k e^{a_k}=a_{\max}+\log\sum_k e^{a_k-a_{\max}}$.

%------------------------------------------------------------------
\section{Analisi del costo sequenziale}
%------------------------------------------------------------------
\begin{center}
\begin{tabular}{lll}
\toprule
Fase & Costo per iterazione & Dati letti \\
\midrule
E-step (distanze, log-sum-exp) & $O(NKD)$ & $X$ ($ND$), $\bth$ ($KD$) \\
Statistiche sufficienti & $O(NKD)$ & $X$, $r$ ($NK$) \\
M-step (formule \eqref{eq:pi}--\eqref{eq:var}) & $O(KD)$ & $S$ ($KD$) \\
Log-likelihood & $O(N)$ & $m_i$ \\
\bottomrule
\end{tabular}
\end{center}
Il costo totale è $T_{\mathrm{seq}}=O(I\cdot NKD)$ con $I$ numero di iterazioni.
Il M-step vero e proprio (aggiornare i parametri a partire dalle statistiche) è
trascurabile: tutto il lavoro è nei due cicli su $N$ punti.

Il rapporto tra dimensioni è il punto chiave: i \textbf{dati} occupano $O(ND)$
(grandi), i \textbf{parametri} $O(KD)$ (piccoli, $K\ll N$).

%------------------------------------------------------------------
\section{Perché EM è parallelizzabile}
%------------------------------------------------------------------

\subsection{Indipendenza dei punti nell'E-step}
Dalla \eqref{eq:estep}, $r_{ik}$ dipende solo da $\bx_i$ e da $\bth^{(t)}$, mai da
altri punti $\bx_j$, $j\ne i$. Una volta noti i parametri, gli $N$ calcoli
delle responsabilità sono \emph{completamente indipendenti} (problema
\emph{embarrassingly parallel}): nessuna dipendenza di dati tra iterazioni del ciclo su $i$.

\subsection{M-step come riduzione}
Le statistiche \eqref{eq:suff} sono somme sui punti di termini che dipendono
ciascuno da un solo $i$. Poiché la somma è associativa e commutativa, si può
partizionare l'insieme dei punti in $P$ blocchi $\mathcal{B}_1,\dots,\mathcal{B}_P$:
\begin{equation}
N_k=\sum_{p=1}^{P}N_k^{[p]},\qquad
S_{kd}^{(1)}=\sum_{p=1}^{P}S_{kd}^{(1)[p]},\qquad
S_{kd}^{(2)}=\sum_{p=1}^{P}S_{kd}^{(2)[p]},
\end{equation}
dove ad esempio $N_k^{[p]}=\sum_{i\in\mathcal{B}_p}r_{ik}$. Ogni blocco calcola
le proprie somme parziali \emph{in locale}, poi una riduzione globale le combina.
Lo stesso vale per la log-likelihood: $\ell=\sum_p\ell^{[p]}$.

\begin{prop}[Correttezza della decomposizione]
Per qualunque partizione $\{\mathcal{B}_p\}$ dell'insieme dei punti, le
statistiche sufficienti globali coincidono (a meno di errori di arrotondamento
in virgola mobile) con quelle calcolate sull'intero dataset.
\end{prop}
\noindent
Questo vale perché la somma è un'operazione associativa e commutativa. Nei numeri
in virgola mobile l'ordine delle somme cambia l'arrotondamento: per questo il
confronto tra versioni va fatto su tolleranza della log-likelihood, non su
uguaglianza bit a bit.

\subsection{I parametri sono piccoli}
Le statistiche sufficienti sono $K(2D+1)$ numeri, indipendenti da $N$. Un
messaggio di $K(2D+1)$ double (più 1 per $\ell$) basta per sincronizzare tutti i
processi. Per esempio, $K=10$, $D=50$ dà $\approx 1010$ double, circa 8 KB.
I dati ($ND$ double) \emph{non si spostano mai}: restano nella memoria del
rank che li possiede.

\subsection{Il M-step finale è così piccolo che si replica}
Dopo la riduzione, ogni rank possiede le stesse statistiche globali. Calcolare
\eqref{eq:pi}--\eqref{eq:var} costa $O(KD)$: più economico che comunicare il
risultato. Ogni rank lo ripete in locale e ottiene gli stessi $\bth^{(t+1)}$,
senza un secondo messaggio.

%------------------------------------------------------------------
\section{Schema di parallelizzazione ibrido}
%------------------------------------------------------------------

\subsection{MPI: partizione dei dati tra rank}
Con $P$ rank, ogni rank $p$ possiede $N_p\approx N/P$ righe di $X$ (distribuzione
a blocchi con \texttt{MPI\_Scatterv}, oppure lettura parallela). I parametri $\bth$ sono
replicati su tutti i rank.

\begin{algorithm}[h]
\caption{Iterazione EM distribuita (rank $p$)}
\begin{algorithmic}[1]
\State \textbf{E-step locale:} per $i\in\mathcal{B}_p$ calcola $r_{ik}$ e $m_i$ \Comment{$O(N_pKD)$}
\State \textbf{Somme locali:} $N_k^{[p]},S^{(1)[p]}_{kd},S^{(2)[p]}_{kd},\ell^{[p]}$ \Comment{$O(N_pKD)$}
\State \texttt{MPI\_Allreduce}(somme locali, \texttt{MPI\_SUM}) $\to$ somme globali \Comment{$O(K D\log P)$}
\State \textbf{M-step replicato:} calcola $\pi,\bmu,\sigma^2$ \Comment{$O(KD)$}
\State Test di convergenza su $\ell$ (identico su tutti i rank)
\end{algorithmic}
\end{algorithm}

Il test di convergenza usa $\ell$ già ridotto con \texttt{MPI\_Allreduce}: tutti i rank
prendono la stessa decisione e non serve broadcast del flag.

\subsection{OpenMP: parallelismo dentro il rank}
Nel ciclo su $i\in\mathcal{B}_p$ ogni thread elabora un sottoinsieme di punti
(\texttt{\#pragma omp parallel for}). Le somme locali sono riduzioni: ogni thread
accumula in un buffer privato, poi i buffer si sommano (clausola
\texttt{reduction} su array, o buffer per thread esplicito). Il rank effettua
l'\texttt{MPI\_Allreduce} una sola volta per iterazione, fuori dalla regione
parallela (o con un solo thread, modello \texttt{MPI\_THREAD\_FUNNELED}).

Perché ibrido? Con $R$ rank per nodo, ciascuno replica $\bth$ e il buffer
dell'\texttt{Allreduce}. Meno rank e più thread riducono la memoria replicata, il
numero di partecipanti alla collettiva e il traffico sulla rete, a parità di core.
Il compromesso rank$\times$thread è uno degli esperimenti previsti.

%------------------------------------------------------------------
\section{Modello di costo e limiti di scalabilità}
%------------------------------------------------------------------
Con $P$ rank, il tempo di una iterazione è circa
\begin{equation}
T_P(N)\;\approx\;\underbrace{\frac{c\,NKD}{P}}_{\text{calcolo}}
+\underbrace{\alpha\log_2 P+\beta\,K(2D+1)\,w\,\log_2 P}_{\text{Allreduce}}
+\underbrace{c'\,KD}_{\text{M-step replicato}},
\label{eq:model}
\end{equation}
dove $c,c'$ sono costi per operazione, $\alpha$ la latenza, $\beta$ l'inverso
della banda e $w$ la dimensione di un double (8 byte). Il modello è per una
\texttt{Allreduce} ad albero; i termini esatti dipendono dall'algoritmo scelto da
MPICH. Dal modello:
\begin{itemize}
  \item il \textbf{calcolo} scala come $1/P$: è la parte che si guadagna;
  \item la \textbf{comunicazione} cresce con $\log P$ e \emph{non dipende da $N$}:
        con $N$ grande è trascurabile, con $N$ piccolo domina;
  \item il \textbf{M-step replicato} è costante in $P$: è la frazione seriale
        (legge di Amdahl), ma piccola perché $O(KD)\ll O(NKD)$.
\end{itemize}

\paragraph{Speedup ed efficienza.}
\begin{equation}
S(P)=\frac{T_1}{T_P},\qquad E(P)=\frac{S(P)}{P}.
\end{equation}
Dalla \eqref{eq:model} l'efficienza resta alta finché $NKD/P\gg
(\alpha+\beta K D w)\log P$, cioè finché ogni rank ha abbastanza punti da
nascondere il costo della collettiva.

\paragraph{Strong vs weak scaling.}
\begin{itemize}
  \item \emph{Strong}: $N$ fisso, $P$ cresce. Il lavoro per rank scende come $1/P$ e prima o poi la
        comunicazione prevale: l'efficienza cala.
  \item \emph{Weak}: $N/P$ fisso. Il calcolo per rank è costante, la comunicazione
        cresce solo come $\log P$: ci si aspetta un'efficienza quasi costante.
\end{itemize}

%------------------------------------------------------------------
\section{Dipendenze di dati e fonti di non-parallelismo}
%------------------------------------------------------------------
Non tutto è parallelo:
\begin{enumerate}
  \item \textbf{Dipendenza tra iterazioni.} L'E-step di $t+1$ richiede $\bth^{(t+1)}$,
        quindi serve una barriera logica (la \texttt{Allreduce}) per iterazione.
        Le iterazioni EM sono inerentemente sequenziali: si parallelizza dentro
        l'iterazione, non tra iterazioni.
  \item \textbf{Convergenza.} Il numero di iterazioni $I$ dipende dai dati e
        dall'inizializzazione, non dal grado di parallelismo.
  \item \textbf{Inizializzazione.} Scelta di $\bth^{(0)}$ (punti casuali o
        $k$-means++) e lettura/distribuzione del dataset: I/O e fase iniziale
        possono essere seriali e vanno misurati a parte.
  \item \textbf{Load imbalance.} Il costo per punto è uguale, quindi una
        partizione a blocchi uniforme è bilanciata. Lo sbilanciamento nasce
        solo da $N\bmod P\neq 0$ o da interferenza di sistema.
  \item \textbf{Memoria.} Il layout di $X$ (riga per riga o per colonna, AoS vs SoA) e l'ordine
        dei cicli $i,k,d$ determinano l'uso della cache. L'intensità
        aritmetica è bassa: il calcolo può essere limitato dalla banda di memoria
        più che dai FLOP.
\end{enumerate}

%------------------------------------------------------------------
\section{Spunti per le ottimizzazioni}
%------------------------------------------------------------------
Dalla teoria seguono direzioni naturali per le tre ottimizzazioni rispetto alla baseline:
\begin{itemize}
  \item \textbf{Fusione dei cicli}: calcolare E-step e accumulo delle statistiche in un solo
        passaggio su $X$, senza memorizzare la matrice $r$ ($N\times K$): meno traffico di memoria.
  \item \textbf{Layout e vettorizzazione}: ordinare i dati e i cicli per accessi contigui, per
        permettere SIMD sul ciclo su $d$.
  \item \textbf{Comunicazione}: impacchettare $N_k,S^{(1)},S^{(2)},\ell$ in un unico buffer,
        una sola \texttt{Allreduce} per iterazione.
  \item \textbf{Ibrido}: scegliere rank$\times$thread per nodo in funzione di $K,D$ e della
        topologia (NUMA).
  \item \textbf{Sovrapposizione} comunicazione/calcolo (\texttt{MPI\_Iallreduce}), utile solo se c'è lavoro
        indipendente da sovrapporre.
\end{itemize}

%------------------------------------------------------------------
\section*{Riferimenti}
%------------------------------------------------------------------
\begin{itemize}
  \item A.~P.~Dempster, N.~M.~Laird, D.~B.~Rubin, \emph{Maximum Likelihood from
        Incomplete Data via the EM Algorithm}, J.\ Royal Statistical Society B, 1977.
  \item C.~M.~Bishop, \emph{Pattern Recognition and Machine Learning}, Springer, 2006, cap.~9.
  \item P.~Pacheco, \emph{An Introduction to Parallel Programming}, Morgan Kaufmann, cap.~3 (MPI).
\end{itemize}

\end{document}
